% ════════════
% WP5_Working_Paper.tex
% Relational Engineering Specifications: Operational Protocols for Outreach Contact and Pod Assignment
% Material Dignity Infrastructure Working Paper 5
% ════════════

\documentclass[12pt, letterpaper]{article}

\usepackage[left=0.7in, right=0.7in, top=1in, bottom=1in, headheight=14pt]{geometry}
\usepackage[expansion=false]{microtype}
\usepackage{textcomp}
\usepackage{setspace}
\usepackage{booktabs}
\usepackage{tabularx}
\usepackage[titles]{tocloft}
\usepackage{tabularray}
\usepackage{longtable}
\usepackage{array}
\usepackage{multirow}
\usepackage{hyperref}
\usepackage{amsmath}
\usepackage{amssymb}
\usepackage{parskip}
\usepackage{titlesec}
\usepackage{fancyhdr}
\usepackage[table]{xcolor}
\usepackage{caption}
\usepackage{footnote}
\usepackage{enumitem}
\usepackage{needspace}
\usepackage{etoolbox}
\usepackage{float}

\definecolor{auditblue}{HTML}{003366}

\setstretch{1.4}
\setlength{\parindent}{0pt}
\setlength{\parskip}{8pt}
\hyphenpenalty=1000

\titleformat{\section}{\normalfont\large\color{auditblue}}{\thesection.}{0.5em}{}
\titleformat{\subsection}{\normalfont\normalsize\color{auditblue}}{\thesubsection.}{0.5em}{}
\titlespacing*{\section}{0pt}{18pt}{6pt}
\titlespacing*{\subsection}{0pt}{12pt}{4pt}

\pagestyle{fancy}
\fancyhf{}
\rhead{\small\textit{{Material Dignity Infrastructure Working Paper 5}}}
\lhead{\small\textit{Working Paper, June 2026}}
\cfoot{\thepage}
\renewcommand{\headrulewidth}{0.4pt}

\renewcommand{\arraystretch}{1.3}

\clubpenalty=10000
\widowpenalty=10000
\displaywidowpenalty=10000
\AtBeginEnvironment{thebibliography}{\interlinepenalty=10000}

\hypersetup{
  colorlinks=true,
  linkcolor=black,
  citecolor=black,
  urlcolor=black,
  pdftitle={{Relational Engineering Specifications}},
  pdfauthor={{Charles J. DiBella}}
}
\setcounter{tocdepth}{2}

\begin{document}
\captionsetup[table]{labelfont=normalfont, labelsep=period, skip=6pt}

\begin{titlepage}
\begin{center}
\setstretch{1.4}
{\LARGE \textbf{Relational Engineering Specifications}}\\[6pt]
{\large Operational Protocols for Outreach Contact and Pod Assignment}\\[0.5cm]
{\normalsize \textbf{Charles J. DiBella}}\\[0.01cm]
{\normalsize Principal Systems Architect}\\[0.01cm]
{\normalsize Working Paper, June 2026}\\[0.3cm]
\textbf{ABSTRACT}
\vspace{0.2cm}
\end{center}
\begin{minipage}{0.95\textwidth}
\setstretch{1.15}
\noindent
Convert the theoretical construct of ``relational dignity'' into a field-executable engineering specification. This paper bridges the ethnographic bedrock of street-level trust formation with the neurobiological architecture of the autonomic nervous system to produce a standardized, scalable operational manual for the Material Dignity Infrastructure (MDI) street-to-home pipeline. Relational dignity is a measurable production output requiring precise structural preconditions. Sustainable transition requires replacing vertical compliance with horizontal, needs-responsive engagement grounded in consistency of return and unconditional regard. This paper operationalizes the Relational Intake Protocol through a three-state Autonomic Routing Matrix, and the Social Network Transition Protocol through a reciprocity-weighted pod assignment algorithm. The protocol's assignment weights and misclassification ceilings are explicit design parameters requiring prospective calibration through the Singular Prototype Threshold, isolating what must be built from what must be proven.
\end{minipage}
\end{titlepage}

\pagenumbering{arabic}

\newpage
\section*{Document Metadata}

\noindent\textbf{Keywords}\\[1pt]
Relational intake protocol, polyvagal theory, autonomic routing, neuroception, social network analysis, pod assignment algorithm, warm offer, street-to-home pipeline, MDI, relational dignity infrastructure, Dunbar pod

\vspace{8pt}
\noindent\textbf{JEL Classification}
\begin{itemize}[nosep, before={\vspace{-8pt}}, leftmargin=2em]
    \item I14. Health and Inequality
    \item I38. Government Policy; Provision and Effects of Welfare Programs
    \item H75. State and Local Government: Health, Education, Welfare
    \item Z13. Economic Sociology; Economic Anthropology; Social and Economic Stratification
\end{itemize}

\vspace{8pt}
\noindent\textbf{Target eJournals}
\begin{itemize}[nosep, before={\vspace{-8pt}}, leftmargin=2em]
    \item Housing and Residential Economics eJournal
    \item Public Health Law and Policy eJournal
    \item Social Capital, Networks \& Trust eJournal
    \item Behavioral \& Experimental Economics eJournal
\end{itemize}

\vspace{8pt}
\noindent\textbf{Suggested Citation}\\[1pt]
DiBella, C.J. (2026). Relational Engineering Specifications: Operational Protocols for Outreach Contact and Pod Assignment. Material Dignity Infrastructure Working Paper 5, June 2026.

\newpage

\vspace{8pt}
\noindent\textbf{Author's Note and Statement of Necessity}\\[1pt]
Working Paper 4 (SSRN 6881539) of this series established that the relational layer must be designed to specification, built into the operational architecture, and measured against defined thresholds. A theoretical framework without operational specifications is a research contribution without deployment capacity. This paper provides the build manual for the two operational decisions that determine whether the relational layer functions from day one of tower operation: who makes the offer and how, and who lives with whom after the offer is accepted. The Autonomic Routing Matrix is the manual for the outreach worker on the street. The Social Network Analysis pipeline is the manual for the Process 4 data administrator managing pod assignments.

The Material Dignity Infrastructure series operates as a cumulative preprint architecture in which each working paper builds on prior papers in the sequence. Citations to prior MDI working papers operate as citations to the developing theoretical framework they establish, not as citations to peer-reviewed empirical findings. Those internal assumptions await external validation through the Singular Prototype Threshold.

\newpage
\begingroup
\setstretch{1.35}
\setlength{\cftbeforesecskip}{2pt}
\setlength{\cftbeforesubsecskip}{-2pt}
\tableofcontents
\endgroup

\clearpage
\setstretch{1.35}

\clearpage
\section{The Ethnographic Bedrock and Theoretical Grounding}

The Autonomic Routing Matrix is a theoretical formalization of Porges's polyvagal architecture, applied to the specific contact conditions of street outreach. Longitudinal street outreach practice on Skid Row, Los Angeles provided the motivating context for this formalization. Field methodology dictates that trust is built through fixed-schedule, non-contingent material offerings (food, tobacco) that signal unconditional regard rather than incentivize behavior. The operative field protocol replaces institutional intake questions with a single inquiry: \textit{do you need anything?}

This foundational field data is partial. It systematically documents successful contact initiation rates and material exchange patterns, but lacks longitudinal tracking of autonomic state shifts and physiological baselines. It serves as the motivating context that demands the protocol, not its empirical validation. The foundation of the protocol relies on the polyvagal literature. The field observations confirm the practical necessity of applying that literature to this population: a fixed-schedule, needs-responsive methodology produced observable affective shifts across repeated encounters with chronically street-present individuals, a pattern consistent with the autonomic state-transition dynamics polyvagal theory predicts.

The critical applied finding is that approach vector, physical posture, and verbal script cannot be fixed across all contact attempts. They must match the observable neurological state of the individual at the moment of contact. Asking ``do you want to come inside?'' to an individual in dorsal-vagal shutdown is a category error because the question requires prefrontal processing capacity that the shutdown state has suspended. The offer fails not because the individual rejects housing but because the offer was made to a nervous system that cannot receive it. Polyvagal theory explains this mechanism. The Autonomic Routing Matrix operationalizes it.

\clearpage
\section{Gap 1: The Relational Intake Protocol}

\subsection{The Neuroception Architecture}

The Autonomic Routing Matrix in the following subsection is grounded in the behavioral taxonomy of polyvagal theory, not in its neurophysiological hierarchy. This distinction is operationally significant. The three-state behavioral taxonomy, which encompasses shutdown, mobilized defense, and social engagement, identifies observable biological indicators that workers can assess in field conditions without resolving the underlying neurophysiological debate about the dorsal vagal complex. Grossman and Taylor (2007), which this paper cites, is not a validation study; it is a partial defense of polyvagal theory that simultaneously raises critiques of specific neurophysiological claims. Polyvagal theory's clinical and applied behavioral framework has broad adoption in trauma-informed care, somatic therapy, and crisis intervention; the underlying neurophysiology remains subject to ongoing empirical debate in the neuroscience literature. The protocol depends on the behavioral distinctions being real and observable, which the clinical literature supports. It does not depend on the resolution of the neurophysiological contest, and it should be read accordingly.

Porges's polyvagal theory identifies neuroception as the nervous system's non-conscious process of continuously scanning the environment for cues of safety or threat, operating below conscious awareness and prior to any behavioral response. Chronic street exposure locks neuroception into one of two defensive configurations. Dorsal-vagal shutdown, the phylogenetically oldest defense state, is characterized by immobility, metabolic conservation, flat affect, and dissociation. Sympathetic arousal, the fight-or-flight state, is characterized by hypervigilance, rapid movement, elevated heart rate, and rejection of proximity.

The clinical distinction between these states matters because the standard institutional intake process treats all contact refusals as equivalent. A person sitting motionless on a milk crate beside the Los Angeles River for twelve hours, eyes unfocused, breathing shallow, body weight collapsed into the seat, is in a qualitatively different neurological condition than a person pacing the perimeter of an encampment, scanning for threats, speaking rapidly, and rejecting the approach of any stranger. The first individual's nervous system has classified the environment as inescapably dangerous and has responded by shutting down voluntary motor activity, including the social engagement system that processes language and evaluates offers. The second individual's nervous system has classified the environment as actively dangerous and has flooded the body with adrenaline to prepare for combat or escape. Both individuals will refuse a housing offer. The refusals mean different things, require different responses, and produce different outcomes when recorded in an institutional database.

A warm offer delivered to either defensive state is structurally invalid. Refusal in dorsal-vagal shutdown is not a cognitive decision. It is an autonomic defense response that the prefrontal cortex did not generate and cannot be attributed to the individual's considered preferences. Refusal in sympathetic arousal is an adrenaline-driven rejection of proximity that will resolve when the physiological cycle completes. Recording either as a housing refusal and closing the application is an institutional category error that converts a neurological state into a policy outcome.

The Relational Intake Protocol begins with state assessment. The worker evaluates observable biological indicators rather than behavior, compliance, or presentation, and routes the contact to the protocol appropriate to the identified state.

\subsection{The Autonomic Routing Matrix}

\textbf{State 1: Dorsal Vagal (Shutdown)}

\textit{Observable indicators:} Total or near-total immobility. Slumped or collapsed posture. Shallow or imperceptible breathing. Flat or absent affect. Vacant gaze. Apparent dissociation from immediate environment. Does not track the approach of the worker.

\textit{Threat perception:} Extreme or inescapable. The organism conserves metabolic energy by suspending voluntary activity.

\textit{Approach vector:} Oblique approach at minimum six feet. Do not approach head-on. Direct sustained eye contact is forbidden because the shutdown nervous system reads it as predation.

\textit{Physical posture:} Lower center of gravity to match the individual's level. Sit or crouch if the individual is on the ground. Hands visible and still.

\textit{Script constraint:} Interrogative questions of any form are forbidden. "How are you?" "Do you want to come inside?" "Can I help you?" All interrogatives demand prefrontal processing capacity that does not exist in State 1.

\textit{Mandatory script:} Declarative statements of safety and non-demand only. "I am going to leave this water here." "I am going to sit here for two minutes, and then I will leave." "I will come back on Thursday."

\textit{Operational routing:} Deliver material supply without requiring acknowledgment. Log the contact, the state, and the material provided. Withdraw. Repeat on the established schedule until autonomic shift to State 2 or State 3 is observed. Do not record as a service refusal.

\textbf{State 2: Sympathetic (Mobilized)}

\textit{Observable indicators:} Pacing or repetitive movement. Rigid musculature. Rapid chest-level breathing. Hyper-vigilant scanning of environment. Accelerated or pressured speech. Sweating.

\textit{Threat perception:} Active and immediate. The organism is flooded with adrenaline and prepared to fight or flee.

\textit{Approach vector:} Maintain distance. Never block an exit path. Position parallel to the individual, facing the same direction rather than face-to-face. This posture signals no predation intent.

\textit{Physical posture:} Standing, relaxed, hands open and visible. Slow all physical movements to approximately half normal speed.

\textit{Script constraint:} Forbidden: commands to calm down, logical arguments about the offer, or any statement that requires the individual to evaluate a choice. Adrenaline forecloses deliberation.

\textit{Mandatory script:} Match the energy level but lower the pitch. Short concrete sentences. Validate the threat perception without confirming delusion. "You are moving. I am staying right here. You have space to move."

\textit{Operational routing:} Allow the adrenaline cycle to complete. This process typically requires five to fifteen minutes of sustained non-escalation. Do not introduce the housing offer until physiological arousal slows, meaning breathing deepens, pacing stops, and scanning frequency decreases. Record the contact and the state.

\textbf{State 3: Ventral Vagal (Engaged)}

\textit{Observable indicators:} Relaxed or open posture. Deep belly breathing. Spontaneous eye contact at normal duration. Tonal variation in voice. Facial micro-expressions of recognition or interest.

\textit{Threat perception:} Absent or low. The organism is capable of social engagement and prefrontal cognitive processing.

\textit{Approach vector:} Direct approach. Face-to-face orientation is appropriate. Normal proximity.

\textit{Physical posture:} Relaxed, mirroring the individual's posture. Normal eye contact.

\textit{Script constraint:} The prefrontal cortex is online. Choices, logistics, and future planning are biologically accessible.

\textit{Mandatory script:} Execute the full warm offer. Reference specific historical knowledge to demonstrate relational continuity. An example of this specificity is "I know you need Buster to come with you. We have a secure kennel on the third floor with his name on it. Your cart is already stored. The room has a lock and the key is yours." Specificity is the evidence of relationship.

\textit{Operational routing:} Proceed to pod transition logistics. The offer is valid.

\subsection{The Intake Failure Rule}

An intake offer made while the individual is in State 1 or State 2 is structurally invalid regardless of what response is recorded. Refusal in State 1 or State 2 is not a refusal of housing. It is an autonomic defense response. The offer must be withheld or withdrawn until State 3 is achieved through sustained co-regulation. Co-regulation is the process by which the worker's own ventral vagal state, maintained through calm and non-demanding presence, gradually shifts the individual's neuroception toward safety.

The co-regulation process is not instantaneous. Field observation on Skid Row documented that individuals in sustained dorsal-vagal shutdown required between three and twelve weekly contacts before autonomic shift became observable. Each contact followed the State 1 protocol. Each contact deposited a small increment of environmental safety evidence into the individual's neuroceptive record. The shift, when it occurred, was observable in specific biological markers. Breathing deepened. Eye tracking resumed. Head position changed from collapsed to upright. These markers preceded verbal engagement by one to three contacts. The worker who recognized the shift could adjust the protocol from State 1 declarative statements to State 3 conversational engagement. The worker who did not recognize the shift continued depositing State 1 contacts into an individual whose nervous system had already moved to State 3 readiness, wasting the window.

Recording a State 1 or State 2 response as a housing refusal and closing the application falsifies the individual's decision-making capacity. The MDI operational protocol prohibits this recording and requires the contact to be logged as a state observation with a scheduled follow-up.

\subsection{State Transition Dynamics}

Autonomic states are not fixed. Individuals cycle between states within a single day and across contact encounters. An individual observed in State 1 on Tuesday may present in State 2 on Thursday because of an intervening encampment sweep, a physical altercation, or a change in substance use pattern. An individual who achieved State 3 during the previous contact may regress to State 1 after a night of rain exposure, theft of possessions, or police contact.

The Routing Matrix requires state assessment at each contact, not a persistent label. The worker does not approach an individual on the basis of the state recorded at the previous visit. The worker approaches on the basis of the state observed in the first thirty seconds of the current encounter. Observable biological indicators are assessed from a distance before verbal contact is initiated. This assessment protocol prevents the most common field error, which is approaching a State 1 individual with the conversational familiarity earned during a previous State 3 encounter. That familiarity triggers threat detection in the shutdown nervous system because the approach vector and energy level are mismatched to the current autonomic reality.

The transition from State 1 to State 3 is not linear. The typical pattern observed in longitudinal outreach contact is State 1 to State 2 to State 1 to State 2 to State 3, with reversion episodes at each environmental disruption. Workers trained in the Routing Matrix expect this cycle and do not interpret reversion as failure. Each contact that matches the protocol to the observed state deposits co-regulation evidence regardless of whether the contact produces a housing acceptance.

\subsection{Minimum Worker Training Specification}

The Autonomic Routing Matrix requires outreach workers to reliably distinguish three distinct autonomic states in field conditions, under time pressure, across individuals with different behavioral baselines and substance use profiles, and without making misclassification errors that cause harm. State misclassification produces the contact failure the protocol was designed to prevent: a State 1 individual approached with State 3 energy receives a threat signal rather than a co-regulation signal, and a State 2 individual presented with an offer requiring deliberation receives a demand the adrenaline cycle cannot process. Both errors damage the relational record at the moment of contact.

The minimum training requirement before a worker operates independently under the Autonomic Routing Matrix is the following. Twenty-four hours of classroom instruction covering polyvagal theory foundations, the three-state behavioral taxonomy, the observable biological indicators for each state, the prohibited and mandatory scripts for each state, and the non-expiring offer protocol. Forty hours of supervised field practice under a certified field supervisor, with a minimum of ten documented State 1 contacts, ten documented State 2 contacts, and ten documented State 3 contacts, each reviewed by the supervisor within forty-eight hours and graded for protocol adherence. A field supervisor must accompany any worker who has not yet completed the forty-hour supervised practice requirement on all contacts involving State 1 or State 2 individuals. These hours are calibrated against two comparable field training standards. Crisis Intervention Team training, the national standard for behavioral crisis response adopted by law enforcement agencies across the United States, requires forty hours of training before independent deployment; the MDI supervised field component matches this floor and applies it to a non-law-enforcement context where the relational consequences of misclassification are qualitative rather than physical. Motivational Interviewing certification under the MINT standard requires fourteen to twenty hours of didactic instruction plus supervised practice; the MDI twenty-four-hour classroom requirement exceeds this floor because the three-state behavioral taxonomy requires more complex discrimination training than the spirit-and-method framework MI instruction covers. The specification intentionally exceeds minimum benchmarks rather than meeting them because outreach workers operating in State 1 and State 2 contact conditions face a lower margin for protocol error than workers in office-based clinical settings.

The acceptable misclassification rate is defined as fewer than one in twenty contacts resulting in a documented protocol breach, meaning a state assessment recorded by the worker and subsequently reviewed and corrected by a supervisor. This threshold is more stringent than comparable behavioral observation standards in adjacent clinical training programs: Crisis Intervention Team (CIT) training, the forty-hour law enforcement crisis response standard, achieves inter-rater agreement on behavioral state classification in the range of eighty to eighty-five percent under structured training conditions, implying a disagreement rate of fifteen to twenty percent. Motivational Interviewing fidelity coding under the MITI standard accepts up to ten to fifteen percent coding disagreement between trained raters. The MDI five-percent ceiling is set more stringently than both comparators because field contact errors produce direct relational harm at the moment of contact.\\
\textbf{Epistemic Calibration:} This five-percent misclassification ceiling is a provisional design parameter requiring prospective calibration at the prototype stage. It is an asserted target, not an empirically validated constant for this workforce. If field data from the Singular Prototype Threshold shows that five percent is unachievable at the training duration specified, the specification will be revised accordingly. Workers exceeding this rate receive additional supervised practice before independent contact authorization is reinstated. Annual recertification requires eight hours of review training and a ten-contact field assessment reviewed by a supervisor. Every contact log entry must record the state assessed at approach, any state revision observed during the contact, and the protocol executed. When a supervisor reviews a contact log and identifies a likely misclassification, the correction is documented and the contact record is amended. The amended record is retained in the resident file and used in aggregate to assess protocol adherence rates across the team.

\subsection{The Non-Expiring Offer Protocol}

The warm offer does not expire. The individual who declines in State 3 receives the same offer at the next scheduled contact. The word "decline" is not used in MDI operational documentation because it implies a considered decision that the protocol cannot assume was made under conditions of full neurological availability. The offer is ongoing. The relationship is the delivery mechanism. The offer becomes credible over time precisely because it continues to arrive without consequence for previous responses.

This protocol departs from standard outreach practice, which typically imposes offer expiration windows ranging from seventy-two hours to thirty days. Expiration windows convert a relational process into a transactional deadline. They communicate that the institution's patience is finite and that the individual's access to housing is contingent on a timeline set by the institution rather than by the individual's neurological readiness. The MDI protocol eliminates this deadline because the field evidence demonstrates that trust accumulates on a biological timeline that institutional calendars cannot compress.


\section{Gap 2: The Social Network Transition Protocol}

\subsection{The Retention Stakes}

Tsai and colleagues' 2012 longitudinal study across three US Housing First sites examined community participation and community integration among formerly homeless adults, finding that social integration was a meaningful correlate of stable housing outcomes. That study does not operationalize social network density as a statistically independent predictor of twelve-month retention in a dose-response relationship; that characterization exceeds what the citation supports and is corrected here. The directional evidence, namely that social connection at ninety days correlates with housing stability, draws on Tsai et al. (2012) and Hawkins and Abrams (2007), whose analysis of network disappearance post-housing transition establishes that social network severance at intake is a retention liability. Neither study produces the specific threshold of three as a measurable floor. The three-contact floor is a design parameter derived from Dunbar's social network architecture: Dunbar (1992, 1998) identifies three to five stable reciprocal relationships as the minimum functional social unit below which individuals lack a reliable co-regulation base, a finding robust across cultural and ecological contexts. That architecture is adapted here to the specific conditions of permanent supportive housing, where the relevant unit is verified reciprocal material-reliance bonds rather than self-reported friendships. The ninety-day Named Peer Contact floor of three is therefore a theoretically grounded design specification requiring empirical calibration at the Singular Prototype Threshold stage, not a figure derived from a single housing stability finding. The prototype will generate the first direct test of whether three is the correct threshold for this population and context.



The dose-response relationship carries a specific implication for tower design. A resident who arrives with two verified reciprocal contacts at intake starts the ninety-day measurement window with a social network that already approaches the Named Peer Contact Index threshold of three. A resident who arrives through lottery assignment starts with zero. The first resident needs to form one new reciprocal bond in ninety days to cross the threshold. The second resident needs to form three. The environmental conditions required to produce three new reciprocal bonds from zero in ninety days inside a building full of strangers are substantially more demanding than the conditions required to produce one additional bond inside a pod where two verified partners already provide social anchoring. The Social Network Transition Protocol exists to guarantee that the maximum number of residents arrive with the relational head start that reduces the gap between intake and threshold.

Encampment communities contain existing social networks with documented reciprocal relational bonds. These networks are invisible to institutional intake systems because institutional intake systems do not ask relational questions. They ask diagnostic questions, compliance questions, and eligibility questions. Housing assignment processes that ignore these networks sever them at the housing threshold and convert a relational asset into a retention liability. Lottery assignment, queue position, and geographic availability each produce this severance. Pod assignment is a clinical decision. The MDI system treats it as one.

\subsection{The Assessment Instrument}

The Social Network Analysis assessment is conducted by the outreach team during the Phase Zero encampment engagement period, prior to tower activation. It uses material reliance questions rather than subjective friendship questions because self-reported friendship is unreliable under conditions of chronic stress and survival competition. Material reliance captures who watches your gear, who you share food with, and who you find first in a crisis. These questions produce a behavioral record that can be independently verified.

The assessment methodology draws on the 1994 foundational Social Network Analysis architecture developed by Wasserman and Faust, adapted to the specific conditions of street-based data collection. Standard instruments assume literate respondents in controlled environments completing written surveys. Street-based assessment requires verbal administration in uncontrolled environments with respondents whose cognitive availability varies by autonomic state. The four-domain interview script was designed to be administerable in under five minutes during a State 3 contact window, using concrete material questions rather than abstract relational ones.

\needspace{10\baselineskip}
\begin{table}[H]
\renewcommand{\arraystretch}{1.2}
\caption*{\textbf{Mandatory Interview Script}}
\begin{tabularx}{\textwidth}{@{} >{\raggedright\arraybackslash}p{3.5cm} >{\raggedright\arraybackslash}X >{\raggedright\arraybackslash}p{5.5cm} @{}}
\toprule
\textbf{Domain} & \textbf{Question} & \textbf{Target Output} \\
\midrule
\textbf{Material Security} & ``When you leave your tent, who watches your gear?'' & Identifies primary trust node \\
\textbf{Metabolic Supply} & ``Who do you share food or tobacco with?'' & Identifies routine reciprocal nodes \\
\textbf{Crisis Response} & ``If you get sick or hurt, who do you find first?'' & Identifies apex reliance node \\
\textbf{Toxicity Flag} & ``Is there anyone here you need to stay away from?'' & Identifies predator or parasite nodes requiring separation \\
\bottomrule
\end{tabularx}
\caption*{Table A: SNA mandatory interview script, four-domain material reliance instrument.}
\end{table}

The sequencing of these questions is deliberate. The first two questions concern routine, low-stakes daily interactions. They establish a conversational pattern of naming people before the third question escalates to crisis-level reliance and the fourth question introduces the sensitive domain of threat assessment. Reversing the sequence, by opening with the toxicity question, would activate threat detection in the respondent and compromise the reliability of the remaining answers.

All responses are logged by name into the Process 4 cybernetic grid. Unnamed responses are discarded.

\subsection{Verification and Scoring}

All claimed ties require bidirectional verification. The worker asks Node A if they watch Node B's gear. Then asks Node B if Node A watches their gear. Unverified claims are algorithmically weighted at zero and excluded from assignment logic. This prevents socially dominant individuals from manufacturing relational claims that inflate their pod placement priority.

\needspace{8\baselineskip}
\begin{table}[H]
\renewcommand{\arraystretch}{1.2}
\begin{tabularx}{\textwidth}{@{} >{\raggedright\arraybackslash}p{3.5cm} >{\raggedright\arraybackslash}X >{\raggedright\arraybackslash}p{5.5cm} @{}}
\toprule
\textbf{Metric} & \textbf{Verification Action} & \textbf{Algorithmic Weight} \\
\midrule
\textbf{Reciprocity, Verified} & Both nodes confirm the tie independently & \textbf{+3}, Primary Cluster Anchor \\
\textbf{Reciprocity, Unverified} & One node denies or is unavailable & \textbf{0}, Discard \\
\textbf{Density, Hub} & Node named by three or more individuals in the encampment & \textbf{+2}, Central Hub \\
\textbf{Toxicity, Risk} & Node named as threat by any verified node & \textbf{$-5$}, Isolation Flag \\
\bottomrule
\end{tabularx}
\caption*{Table B: Tie verification protocol and provisional algorithmic weighting schema.}
\end{table}

\textbf{Epistemic Calibration:} The algorithmic weights assigned to verified ties (+3 for reciprocity, +2 for hub density, -5 for toxicity risk) are provisional pilot parameters. They are not derived from empirical calibration of homelessness network data. They represent a theoretical weighting of the relative impact of each tie type on housing retention. These weights require explicit calibration against real-world retention data during the Singular Prototype Threshold deployment.

\subsection{Process 4 Assignment Algorithm}

The scored data enters the Process 4 cybernetic grid. Pod assignment executes four constraints in strict priority order.

\textbf{Priority One, Isolation Constraint.} Any node carrying a Toxicity Flag at negative five is assigned to a physically separate floor or building from every node that named them as a threat. This constraint overrides all others. No exception.

\textbf{Priority Two, Anchor Constraint.} All mutually verified Reciprocity pairs at positive three are mandatorily assigned to the same twelve-person Dunbar sub-pod within their floor. Verified bond partners who cannot be co-assigned within capacity constraints trigger a capacity expansion review before any separation is recorded.

\textbf{Priority Three, Hub Constraint.} High-Density nodes at positive two are assigned to pods identified as requiring social anchors. These pods contain a high proportion of individuals carrying zero verified ties. Placing high-density nodes in these environments deploys their existing social gravity to serve the pod's community formation rather than concentrating social capital in already-dense clusters.

\textbf{Priority Four, Remnant Constraint.} Individuals with zero verified ties are distributed evenly across high-density pods. No pod may contain a majority of zero-tie individuals. Social isolation is a contagious environmental condition. Concentrating isolated individuals produces a pod environment in which no community formation is possible regardless of physical design quality.

\subsection{The Assignment Failure Rule}

Assigning mutually verified bond partners at positive three to separate floors constitutes an architectural failure. The social network transition protocol identifies it as equivalent to offering a housing unit without a lock. The material infrastructure is present but the relational layer has been destroyed at intake.

The cascade from this failure is predictable. Separated bond partners lose their primary co-regulation resource at the moment of maximum environmental stress. The housing transition itself is a high-stress event because the individual moves from a known environment with established threat patterns into an unknown environment with unestablished threat patterns. The verified bond partner is the one element of the new environment that the individual's nervous system has already classified as safe. Removing that element forces the individual's neuroception to classify the entire new environment as unknown, which produces the defensive autonomic states that the Relational Intake Protocol was designed to prevent at the street level. The separation recreates the intake problem inside the building.

A pod assignment log that separates verified bond partners without a documented capacity justification and a documented escalation path for review is an operational violation of the RDI retention standard. The Process 4 data administrator who executes such a separation is required to file a capacity exception report that triggers review by the Pod Steward and the tower operations lead within forty-eight hours. The review determines whether the separation can be reversed through a floor-level reassignment or whether the capacity constraint is genuine and permanent.

\subsection{Process 4 Conflict Resolution Protocol}

The four-priority algorithm resolves clear cases. Three conflict configurations occur in real encampment social networks that the algorithm cannot resolve without a human override, and the override authority and decision criteria must be specified before the algorithm is operable as a field tool.

\textbf{Conflict Type A: Competing verified bond pairs for the same pod slot.} When two verified reciprocal pairs are both assigned to Priority Two placement and the available pod has capacity for only one pair, the Pod Steward holds override authority. The decision criterion is recency of verification: the pair whose bidirectional verification was completed most recently receives the slot, on the grounds that the verification reflects the current relational state rather than a historical one. The displaced pair triggers a capacity expansion review filed within twenty-four hours with the Tower Operations Lead.

\textbf{Conflict Type B: Hub node assigned pod is at capacity.} When a Priority Three hub node cannot be placed in an identified anchor pod because that pod has reached its twelve-person ceiling, the Tower Operations Lead holds override authority. The decision criterion is network reach: the hub node is reassigned to the pod whose current residents have the lowest aggregate verified-tie count, so the hub's social gravity addresses the greatest integration deficit. The reassignment is documented with the original assignment record retained.

\textbf{Conflict Type C: Toxicity flag conflicts with a strong reciprocity tie between flagged and non-flagged individuals.} When a node carries a Priority One isolation flag and also holds a verified positive-three reciprocity tie with a non-flagged node, the two constraints are irreconcilable within the standard algorithm. The MDI Regional Director holds override authority for this conflict type, given the combination of safety risk and relational severance that any resolution creates. The decision criteria are applied in this order: first, whether the toxicity flag was filed by the individual who holds the reciprocal tie with the flagged node, in which case the flag overrides the tie without exception; second, whether the toxicity flag was filed by a third party and the reciprocal-tie holder disputes it, in which case a review interview is conducted with both the filer and the tie holder before the final assignment is executed. All Conflict Type C resolutions are documented in a permanent flag-review file available for audit.

All override decisions, regardless of type, must be documented in the Process 4 assignment log with the conflict type identified, the overriding authority named, the decision criterion applied, and the outcome recorded. No verbal override authorizations are accepted. The assignment log is a legal document for the purposes of capacity exception and RDI compliance review.

\clearpage
\section{Integration with WP4 Verification Metrics}

The five RDI leading indicators established in WP4 acquire operational meaning from the parameters in this paper. Three of the five trace to the protocols above.

The \textbf{Named Peer Contact Index} requires three or more reciprocal contacts per resident by day 90. It is a direct output of the Social Network Transition Protocol. Pods assembled through the Anchor and Hub constraints arrive with pre-existing verified ties, accelerating the progress toward the three-contact threshold. A pod assembled through lottery arrives with zero verified ties and must build from nothing inside the tower environment.

The \textbf{Voluntary Service Engagement Rate} requires seventy percent engagement or higher by day 90. It is a direct output of the Relational Intake Protocol. Residents who entered through a State 3 warm offer delivered by a known person with specific historical knowledge arrive with an existing institutional trust foundation. Residents who entered through a State 1 or State 2 offer that was improperly validated arrive with institutional trust at zero or below zero and must rebuild from a trust deficit.

The \textbf{Sustained Ground Floor Opt-Out Rate} must remain at fifteen percent or below over 90 days. It is partially traceable to pod assignment quality. Residents with verified relational bonds assigned to their pod have social reasons to leave their unit and enter common spaces. Isolated residents assigned to pods where they know no one have no social incentive for common space engagement and are precisely the population the fifteen percent threshold is designed to detect as a system signal.

The fiscal calculus from WP4 Section VII is reproduced here with its derivation made explicit. All per-resident figures in the table represent ninety-day avoided-cost fractions, aligned to the Named Peer Contact Index measurement window, rather than annualized projections. The annual baseline figures from which the ninety-day fractions are derived are stated in the source column.

\begin{table}[H]
\renewcommand{\arraystretch}{1.2}
\begin{tabularx}{\textwidth}{@{} >{\raggedright\arraybackslash}p{5cm} >{\raggedright\arraybackslash}X >{\raggedright\arraybackslash}p{3.5cm} @{}}
\toprule
\textbf{Variable} & \textbf{Source and Assumption} & \textbf{Value} \\
\midrule
\textbf{Tower pods} & MDI architecture: 13 Dunbar sub-pods per tower & 13 \\
\textbf{Residents per pod} & MDI pod specification: 12 per sub-pod & 12 \\
\textbf{Total residents} & 13 pods $\times$ 12 residents & 156 \\
\textbf{Residents above Named Peer Contact threshold} & Conservative estimate: 50\%--67\% of residents reach the 3-contact floor through the protocol & 78--104 \\
\textbf{Annual avoided cost per resident (primary)} & California State Auditor (2024): \$50,000/year per-person street-management cost, Los Angeles County, inclusive of ED, law enforcement, and sanitation & \$50,000/year \\
\textbf{Annual avoided cost per resident (historical baseline)} & Culhane et al.\ (2002): \$16,282/year in avoided public service costs for Housing First participants. Retained for trend context; 2002 figures predate LA cost inflation & \$16,282/year \\
\textbf{90-day avoided cost, low estimate} & 78 residents $\times$ (\$16,282 $\div$ 4) & \$317,000 \\
\textbf{90-day avoided cost, primary estimate} & 78--104 residents $\times$ (\$50,000 $\div$ 4) & \$975,000--\$1,300,000 \\
\bottomrule
\end{tabularx}
\caption*{Table C: Fiscal derivation, Named Peer Contact floor avoided-cost range (90-day window), before Medi-Cal billing.}
\end{table}

The primary estimate uses the California State Auditor Los Angeles County figure as the avoided-cost anchor because it reflects current street-management unit costs and is jurisdiction-specific. The Culhane historical baseline is retained to show trend direction; the gap between \$16,282 in 2002 and \$50,000 in 2024 reflects twenty-two years of emergency department, law enforcement, and sanitation cost inflation in Los Angeles County. These figures represent avoided downstream processing costs generated by residents who cross the Named Peer Contact threshold and sustain housing stability. The Social Network Transition Protocol is the operational mechanism that produces the threshold-crossing rate. The warm offer relational standard determines whether the resident arrives capable of achieving it.




\clearpage

\clearpage
\section{Contribution and Sequence}

This paper makes three contributions to the MDI-RDI series.

It converts theoretical production conditions into operational field protocols. WP4 established that the warm offer must be relational from its first moment and that social network bonds must be preserved across the housing transition. This paper details how both requirements are operationally satisfied at the level of approach vector, verbal script, scoring algorithm, and assignment constraint. Strict protocol parameters are what make replication possible. Without them, each implementation team improvises, and improvisation produces the variable quality that defeats the Singular Prototype Threshold's replication logic.

It applies polyvagal theory's three-state architecture as an operational field protocol for a population the clinical literature has not previously addressed in this form. The Autonomic Routing Matrix derives its behavioral parameters from the polyvagal literature. Street outreach practice on Skid Row provided the applied context that confirmed the theory's relevance to chronic street exposure and motivated its formalization into protocol form. The biological state of the individual at the moment of contact determines whether the contact produces trust accumulation or trust damage, regardless of the skill or goodwill of the worker. This is a polyvagal prediction, and the protocol is designed to operationalize it.




\newpage


\clearpage
\addcontentsline{toc}{section}{References}
\section*{References}

Community Guide. "Housing First: Community Living for Adults Experiencing Homelessness." \textit{The Community Preventive Services Task Force}, 2020.

Culhane, Dennis P., Stephen Metraux, and Trevor Hadley. "Public Service Reductions Associated with Placement of Homeless Persons with Severe Mental Illness in Supportive Housing." \textit{Housing Policy Debate} 13, no. 1 (2002): 107–163.

DiBella, Charles J. \textit{Material Dignity Infrastructure: A Distributed Stewardship Model for Homelessness Resolution}. SSRN Working Paper 5968756. 2025.

DiBella, Charles J. \textit{Material Dignity Infrastructure: Structural Misalignment and the Activation of Surplus Shelter Capacity}. SSRN Working Paper 6211658. February 2026.

DiBella, Charles J. \textit{Material Dignity Infrastructure: Los Angeles Metropolitan Stabilization: A Street-to-Home Pipeline Analysis}. SSRN Working Paper 6579600. May 2026.

DiBella, Charles J. \textit{Relational Dignity Infrastructure: The Human Layer Required to Make Material Housing Inhabitable}. SSRN Working Paper 6881539. June 2026.

Dunbar, Robin I.M. "Neocortex Size as a Constraint on Group Size in Primates." \textit{Journal of Human Evolution} 22, no. 6 (1992): 469--493.

Dunbar, Robin I.M. \textit{Grooming, Gossip, and the Evolution of Language}. Cambridge: Harvard University Press, 1998.

Grossman, Peter, and Edwin W. Taylor. "A Clarification and Defense of the Theory of Polyvagal Control." \textit{Biological Psychology} 74, no. 2 (2007): 260–268.

Hawkins, Robert L., and Craig Abrams. "Disappearing Acts: The Social Networks of Formerly Homeless Individuals with Co-occurring Disorders." \textit{Social Science and Medicine} 65, no. 10 (2007): 2031–2042.

Pleace, Nicholas, Riitta Culhane, Riitta Granfelt, and Marcus Knutagård. \textit{The Finnish Homelessness Strategy: An International Review}. Helsinki: Ministry of the Environment, 2015.

Porges, Stephen W. \textit{The Polyvagal Theory: Neurophysiological Foundations of Emotions, Attachment, Communication, and Self-Regulation}. New York: Norton, 2011.

Tsai, Jack, Richard Stanhope, Robert A. Rosenheck, and Wesley J. Kasprow. "The Role of Housing, Substance Abuse, and Mental Illness in Community Participation and Community Integration." \textit{Psychiatric Services} 63, no. 5 (2012): 435–442.

Wasserman, Stanley, and Katherine Faust. \textit{Social Network Analysis: Methods and Applications}. Cambridge: Cambridge University Press, 1994.

Y-Foundation. \textit{A Home of Your Own: Housing First and Ending Homelessness in Finland}. Helsinki: Y-Foundation, 2022.

\end{document}
